\documentclass[tikz,border=10pt]{standalone}

\usepackage{amsmath,amssymb}
\usepackage{fontspec}
\usepackage{xeCJK}
\IfFontExistsTF{SimHei}{\setCJKmainfont{SimHei}}{\setCJKmainfont{Noto Sans CJK SC}}
\usepackage{xcolor}
\usepackage{tikz}
\usetikzlibrary{arrows.meta,backgrounds,calc,fit,positioning}

\definecolor{ink}{HTML}{18324A}
\definecolor{muted}{HTML}{5F7286}
\definecolor{linec}{HTML}{CBD8E5}
\definecolor{cyanfill}{HTML}{EEF9FB}
\definecolor{cyanedge}{HTML}{6AB7C9}
\definecolor{purplefill}{HTML}{F4F1FB}
\definecolor{purpleedge}{HTML}{BDA8DE}
\definecolor{greenfill}{HTML}{EFFAF4}
\definecolor{greenedge}{HTML}{A8D7B6}
\definecolor{orangefill}{HTML}{FFF6ED}
\definecolor{orangeedge}{HTML}{F2C892}
\definecolor{redfill}{HTML}{FFF1F4}
\definecolor{rededge}{HTML}{E8A7B7}
\definecolor{badgecyan}{HTML}{1F7A8C}
\definecolor{badgepurple}{HTML}{6A4C93}
\definecolor{badgegreen}{HTML}{1B9E77}
\definecolor{badgeorange}{HTML}{D97925}
\definecolor{badgered}{HTML}{C94B59}

\tikzset{
  statebox/.style={
    rectangle,
    rounded corners=10pt,
    draw=#1,
    fill=#1!12,
    minimum width=3.2cm,
    minimum height=1.8cm,
    text width=3.2cm,
    align=center,
    font=\sffamily\small\bfseries,
    inner sep=7pt
  },
  statebox/.default=cyanedge,
  statedesc/.style={font=\sffamily\scriptsize, text=muted, align=center},
  flow/.style={-{Stealth[length=3.4mm]}, line width=0.95pt, draw=ink},
  obsline/.style={line width=0.8pt, draw=muted, dashed},
  labelbox/.style={
    fill=white,
    rounded corners=7pt,
    inner sep=3.5pt,
    font=\sffamily\tiny,
    text=ink
  },
}

\begin{document}
\begin{tikzpicture}[node distance=1.55cm and 2.1cm]

\node[
  rectangle,
  rounded corners=10pt,
  fill=cyanfill,
  inner sep=6pt,
  text=badgecyan,
  font=\sffamily\scriptsize\bfseries
] (title) at (0,0) {OpenAMP 控制状态机};
\node[statedesc, text width=12.4cm, below=0.55cm of title] (subtitle)
{五个核心状态用于约束 JOB\_REQ、JOB\_ACK、HEARTBEAT、SAFE\_STOP 与 JOB\_DONE 的合法转移；STATUS\_REQ/RESP 只承担观测职责，不参与主状态推进。};

\node[statebox=cyanedge, fill=cyanfill, below=1.1cm of subtitle] (ready) {READY};
\node[statedesc, below=0.1cm of ready] {空闲待命\\可接受 JOB\_REQ};

\node[statebox=purpleedge, fill=purplefill, right=2.05cm of ready] (checking) {CHECKING};
\node[statedesc, below=0.1cm of checking] {工件与参数校验};

\node[statebox=greenedge, fill=greenfill, right=2.05cm of checking] (running) {RUNNING};
\node[statedesc, below=0.1cm of running] {执行中 / 心跳监护激活};

\node[statebox=rededge, fill=redfill, below=2.55cm of checking] (fault) {FAULT};
\node[statedesc, below=0.1cm of fault] {故障信息保留\\仅供查询};

\node[statebox=orangeedge, fill=orangefill, below=2.55cm of running] (safestop) {SAFE\_STOP};
\node[statedesc, below=0.1cm of safestop] {安全停机与收敛};

\draw[flow] (ready) -- (checking)
  node[midway, above=4pt, labelbox] {JOB\_REQ};
\draw[flow] (checking) -- (running)
  node[midway, above=4pt, labelbox] {JOB\_ACK(ALLOW)};
\draw[flow] (checking.west) to[out=215,in=325] node[midway, below=6pt, labelbox] {JOB\_ACK(DENY)} (ready.south east);
\draw[flow] (running.east) to[out=35,in=145] node[midway, above=6pt, labelbox] {JOB\_DONE(success)} (ready.east);
\draw[flow] (running) -- (safestop)
  node[midway, right=6pt, labelbox] {SAFE\_STOP / heartbeat timeout};
\draw[flow] (safestop.west) to[out=190,in=350] node[midway, below=8pt, labelbox] {故障记录落盘后恢复 READY} (ready.south);

\draw[obsline] (safestop.west) -- (fault.east)
  node[midway, above=4pt, labelbox] {last\_fault\_code 可观测};

\node[
  rectangle,
  rounded corners=10pt,
  draw=linec,
  fill=white,
  text width=11.4cm,
  align=left,
  inner sep=9pt,
  font=\sffamily\scriptsize,
  text=muted,
  below=1.45cm of fault.south east,
  xshift=2.7cm
] (observe)
{\textbf{STATUS\_REQ / STATUS\_RESP 观测带。} Linux 主核可主动查询 guard\_state、active\_job\_id、heartbeat\_ok、last\_fault\_code 与 total\_fault\_count；该查询不推动状态迁移，因此被单独放到图下方说明。};

\begin{scope}[on background layer]
  \node[
    fit=(title)(subtitle)(ready)(checking)(running)(fault)(safestop)(observe),
    fill=cyanfill!40,
    draw=linec,
    rounded corners=12pt,
    inner sep=18pt
  ] {};
\end{scope}

\end{tikzpicture}
\end{document}
